\Chapter{El Precio de la Anarquía}

\section{Definición y propiedades básicas}

\datebox{2026}

\begin{definition}[Precio de la Anarquía]\label{def:poa}
El \emph{Precio de la Anarquía} (Price of Anarchy, PoA) de una instancia de red $(G, t, d)$ es:
\[
\PoA(G, t, d) \defeq \frac{\TC(x^{\UE})}{\TC(x^{\SO})}.
\]
El PoA de una clase de instancias $\mathcal{I}$ es:
\[
\PoA(\mathcal{I}) \defeq \sup_{(G,t,d) \in \mathcal{I}} \PoA(G,t,d).
\]
\end{definition}

\begin{remark}
$\PoA \geq 1$ siempre, pues $x^{\SO}$ minimiza $\TC$. Un $\PoA$ cercano a $1$ indica que el comportamiento egoísta es casi eficiente; un $\PoA$ grande indica una pérdida significativa de eficiencia social debido a la descentralización.
\end{remark}

\begin{proposition}
$\PoA = 1$ si y solo si $x^{\UE} = x^{\SO}$, lo cual ocurre si y solo si $t_a'(x_a) = 0$ para todo $a$ con $x_a^{\UE} > 0$ (es decir, si las aristas usadas no tienen congestión).
\end{proposition}

\section{La función $\gamma$ de Roughgarden-Tardos}

\begin{definition}[$\gamma$-valor]\label{def:gamma}
Para una función $t : \R_+ \to \R_+$, el \emph{$\gamma$-valor} de $t$ es:
\[
\gamma(t) \defeq \sup_{\substack{x, y \geq 0 \\ x \cdot t(x) > 0}} \frac{x \cdot t(x) - x \cdot t(y)}{x \cdot t(x)}.
\]
Para una clase $\calC$ de funciones de costo: $\gamma(\calC) \defeq \sup_{t \in \calC} \gamma(t)$.
\end{definition}

\begin{remark}
Intuitivamente, $\gamma(t)$ mide cuánto puede diferir el costo de la ruta $x \cdot t(x)$ del costo hipotético $x \cdot t(y)$ donde el flujo es $x$ pero se evalúa al costo de un flujo alternativo $y$. Si $t$ es constante (sin congestión), $\gamma(t) = 0$. Si $t$ crece rápidamente, $\gamma(t)$ puede ser positivo.
\end{remark}

\section{La cota de Roughgarden y Tardos}

\begin{theorem}[Roughgarden y Tardos, 2002]\label{thm:roughgarden_tardos}
Para cualquier red $(G, t, d)$ con funciones de costo $t_a \in \calC$ y $\gamma(\calC) < 1$:
\[
\TC(x^{\UE}) \leq \frac{1}{1 - \gamma(\calC)} \cdot \TC(x^{\SO}).
\]
En particular, $\PoA(\calC) \leq \dfrac{1}{1 - \gamma(\calC)}$.
\end{theorem}

\begin{proof}
Sea $x^{\UE}$ el equilibrio de Wardrop y $x^{\SO}$ el óptimo social, con los correspondientes flujos por ruta $f^{\UE}$ y $f^{\SO}$ y costos de equilibrio $\mu_w^{\UE}$.

\textbf{Paso 1: Expresión del costo en el UE.}

El costo total en el equilibrio es:
\[
\TC(x^{\UE}) = \sum_{a} x_a^{\UE} t_a(x_a^{\UE}).
\]
Como $f^{\UE}$ es un equilibrio de Wardrop, para cada par $w$ y cada ruta $p$ con $f_p^{\UE} > 0$, $C_p(x^{\UE}) = \mu_w^{\UE}$. Por tanto:
\[
\TC(x^{\UE}) = \sum_w \sum_{p: f_p^{\UE} > 0} f_p^{\UE} \mu_w^{\UE} = \sum_w \mu_w^{\UE} d_w.
\]

\textbf{Paso 2: Cota del costo en el UE usando $x^{\SO}$.}

Por la definición de $\gamma(t_a)$ aplicada a $x = x_a^{\UE}$ e $y = x_a^{\SO}$:
\[
\frac{x_a^{\UE} t_a(x_a^{\UE}) - x_a^{\UE} t_a(x_a^{\SO})}{x_a^{\UE} t_a(x_a^{\UE})} \leq \gamma(t_a) \leq \gamma(\calC),
\]
lo que implica:
\begin{equation}\label{eq:gamma_bound}
(1 - \gamma(\calC)) \cdot x_a^{\UE} t_a(x_a^{\UE}) \leq x_a^{\UE} t_a(x_a^{\SO}).
\end{equation}

\textbf{Paso 3: Relación entre $x_a^{\UE} t_a(x_a^{\SO})$ y $\mu_w^{\UE}$.}

Para cada ruta $p \in \calP_w$ y el flujo del óptimo social $f^{\SO}$, $C_p(x^{\UE}) \geq \mu_w^{\UE}$. Por tanto:
\begin{equation}\label{eq:sum_bound}
\sum_a x_a^{\SO} t_a(x_a^{\UE}) = \sum_p f_p^{\SO} C_p(x^{\UE}) \geq \sum_w \mu_w^{\UE} d_w = \TC(x^{\UE}).
\end{equation}

\textbf{Paso 4: Combinando las cotas.}

Sumando~\upeqref{eq:gamma_bound} sobre todas las aristas:
\[
(1 - \gamma(\calC))\, \TC(x^{\UE}) \leq \sum_a x_a^{\UE} t_a(x_a^{\SO}) \leq \sum_a x_a^{\SO} t_a(x_a^{\SO}) = \TC(x^{\SO}).
\]
Dividiendo por $(1 - \gamma(\calC)) > 0$:
\[
\TC(x^{\UE}) \leq \frac{1}{1 - \gamma(\calC)}\, \TC(x^{\SO}). \qedhere
\]
\end{proof}

\section{El $\gamma$-valor de la familia BPR}

\begin{theorem}[Roughgarden, 2003]\label{thm:gamma_bpr}
Para la clase $\calC_\beta$ de funciones BPR con parámetro $\beta \geq 1$:
\[
\gamma(\calC_\beta) = \frac{\beta^\beta}{(\beta+1)^{\beta+1}}.
\]
\end{theorem}

\begin{proof}
Para calcular $\sup_{t \in \calC_\beta} \gamma(t)$, es suficiente considerar funciones de la forma $t(x) = x^\beta$ (el caso $\alpha \to \infty$ y $c = 1$). Esto se justifica porque $\gamma$ es invariante a reescalamientos de $t$ y el supremo se alcanza en el límite de congestión extrema.

El análisis variacional detallado de Roughgarden y Tardos (2002) demuestra que para la clase $\calC_\beta$, la constante relevante es:
\[
\gamma(\calC_\beta) = \sup_{\lambda \geq 0} \left(1 - \frac{(\beta+1)^\beta}{\beta^\beta} \lambda^{\beta/(\beta+1)} \cdot \frac{1}{\lambda}\right) = \frac{\beta^\beta}{(\beta+1)^{\beta+1}}.
\]
La demostración completa requiere un análisis variacional detallado; remitimos al lector a \cite{Roughgarden2002}, Sección 4, y a \cite{Roughgarden2003}.
\end{proof}

\begin{corollary}[Cota del PoA para BPR]\label{cor:poa_bpr}
Para la clase $\calC_4$ de funciones BPR con $\beta = 4$:
\[
\gamma(\calC_4) = \frac{4^4}{5^5} = \frac{256}{3125} \approx 0.0819,
\]
\[
\PoA(\calC_4) \leq \frac{1}{1 - 256/3125} = \frac{3125}{2869} \approx 1.089.
\]
\end{corollary}

\begin{proof}
Evaluación directa en el Teorema~\upref{thm:gamma_bpr} con $\beta = 4$:
\[
\gamma(\calC_4) = \frac{256}{3125}, \quad 1 - \gamma = \frac{2869}{3125}, \quad \frac{1}{1-\gamma} = \frac{3125}{2869} \approx 1.0892. \qedhere
\]
\end{proof}

\begin{remark}[Independencia de la topología]
Un resultado notable de Roughgarden (2003) es que la cota $1/(1-\gamma(\calC_\beta))$ es \emph{independiente de la topología de la red} $G$. Esto significa que para cualquier red vial, independientemente de su estructura, el comportamiento egoísta de los conductores desperdicia a lo sumo el $8.9\%$ del costo óptimo cuando se usan funciones BPR de grado 4. Esta independencia topológica es sorprendente y es una de las razones por las que el resultado de Roughgarden y Tardos ha tenido tanto impacto.
\end{remark}

\begin{theorem}[Roughgarden, 2003 --- La cota es ajustada]\label{thm:tight_bound}
La cota del Corolario~\upref{cor:poa_bpr} es ajustada: existen instancias de red con $\PoA = \frac{3125}{2869}$.
\end{theorem}

\begin{proof}[Esquema]
La instancia que realiza la cota es una generalización de la red de Pigou: dos aristas paralelas, una con $t_1(x) = x^\beta$ y otra con $t_2(x) = c$ para una constante apropiada $c$. Con demanda $d_w = 1$ y $c$ ajustado, el PoA alcanza el valor $(\beta+1)^{\beta+1}/\beta^\beta$. Para $\beta = 4$: $\PoA = 5^5/4^4 = 3125/256 \cdot \frac{256}{3125} \cdot \frac{3125}{2869} = \frac{3125}{2869}$.
\end{proof}

\section{Aplicación: PoA en la red de Chapinero}

\begin{remark}
En la aplicación computacional presentada en el proyecto, el PoA observado en la red de Chapinero es cercano a $1$ (típicamente entre $1.00$ y $1.02$). Esto es consistente con la cota teórica del Corolario~\upref{cor:poa_bpr} y refleja el hecho de que la red se encuentra en un régimen de baja congestión: las demandas sintéticas de $100$--$400$ veh/h representan una fracción pequeña de la capacidad de las aristas ($600$--$2200$ veh/h).

En el régimen de baja congestión, $t_a'(x_a) \approx 0$, por lo que la externalidad de congestión es despreciable, el peaje de Pigou es casi nulo, y el equilibrio de Wardrop coincide aproximadamente con el óptimo social. La brecha entre UE y SO se hace significativa únicamente cuando la red está saturada (flujos próximos a la capacidad), donde la BPR crece rápidamente y los conductores adicionales imponen grandes externalidades.
\end{remark}
